\chapter{System Design}
\label{chap:design}

\section{High-Level Architecture}
\label{sec:architecture}

SHOOTRZ follows a client--server architecture with a clear separation between the mobile frontend and the backend API service, connected through a managed cloud database and storage layer. Figure~\ref{fig:architecture} presents an overview of the system components and their relationships.

\begin{figure}[H]
  \centering
  \figurebox[0.86\textwidth]{images/system_architecture.png}
  \caption{SHOOTRZ high-level system architecture.}
  \label{fig:architecture}
\end{figure}

The three principal layers are:

\begin{itemize}
  \item \textbf{Mobile Frontend}: A React Native/Expo application running on iOS and Android, responsible for video capture and upload, result display, conversational chat, drill browsing, and progress visualisation.
  \item \textbf{FastAPI Backend}: A Python 3.11+ service that handles video analysis, LLM enrichment, job management, and data persistence orchestration. CPU-intensive pipeline work runs inside a \lstinline|ProcessPoolExecutor|, with concurrency limited by an asyncio semaphore.
  \item \textbf{Supabase Platform}: Managed PostgreSQL with Row-Level Security, Supabase Auth (JWT/PKCE), and a Storage bucket for video artefacts. All durable user state lives here.
\end{itemize}

The backend is stateless with respect to HTTP requests; all durable state lives in Supabase. The exception is an ephemeral SQLite job store that tracks in-flight analysis jobs with a 72-hour TTL, allowing the backend to survive restarts within the retention window without losing committed results.

\section{Frontend Architecture}
\label{sec:frontend_arch}

\subsection{Technology Stack}

The mobile application is built with React Native 0.81.5, TypeScript~5.9, and Expo~54. React~19.1 is used as the rendering foundation. Key library choices are listed in Table~\ref{tab:frontend_stack}.

\begin{table}[H]
  \centering
  \caption{Frontend technology stack.}
  \label{tab:frontend_stack}
  \rowcolors{2}{gray!8}{white}
  \begin{tabularx}{\textwidth}{lX}
    \toprule
    \textbf{Library / Tool} & \textbf{Purpose} \\
    \midrule
    React Navigation 7 & Bottom-tab navigator; deep link handling \\
    Axios 1.12 & HTTP client; JWT injection in headers \\
    @supabase/supabase-js 2.76 & Authentication client; anon key JWT verification \\
    expo-camera 17 & In-app video recording \\
    expo-file-system 19 & Artefact download to device storage \\
    react-native-chart-kit & Angle timeline chart (AngleGraph component) \\
    AsyncStorage & Local offline cache (capped at 200 entries) \\
    \bottomrule
  \end{tabularx}
\end{table}

\subsection{State Management}

Application state is managed through three React Context providers:

\begin{itemize}
  \item \textbf{AuthContext}: Manages the Supabase session lifecycle, PKCE OAuth flow, deep link parsing for OAuth redirects, and the new-user detection gate that directs first-time users through onboarding. The PKCE code verifier is generated locally using \lstinline|expo-crypto|; the authorisation code is extracted from the deep link redirect URI \lstinline|shootrz://...?code=...| using a regex-based fallback parser.
  \item \textbf{HistoryContext}: Fetches and caches the user's analysis history from the API, triggering refreshes when the app returns to foreground or after a new analysis is committed.
  \item \textbf{ProfileContext}: Holds the user's extended profile, skill level, and preferences, exposing update methods that write through to Supabase.
\end{itemize}

\subsection{Key UI Components}

A purpose-built design system (implemented in \lstinline|src/theme/tokens.ts|) defines the full colour palette, type scale, spacing units, and motion curves as named tokens. Analysis-specific components include:

\begin{itemize}
  \item \textbf{ScoreRing}: An animated 0--100 dial with tier colour-coding (elite = gold, great = green, good = blue, fair = orange, poor = red).
  \item \textbf{MetricCard}: Displays a single metric result with name, value, verdict, confidence level, and explanation text.
  \item \textbf{AngleGraph}: A frame-indexed angle timeline rendered using react-native-chart-kit.
  \item \textbf{AnalysisOverlayVideo}: An expo-video player for the annotated skeleton overlay video downloaded from Supabase Storage.
\end{itemize}

\section{Backend Architecture}
\label{sec:backend_arch}

\subsection{Framework and Entry Point}

The backend is implemented using FastAPI 0.110+ on Uvicorn. The application is initialised inside an \lstinline|@asynccontextmanager| lifespan function that enforces \lstinline|multiprocessing.set_start_method("spawn")| on non-Windows platforms---a requirement because MediaPipe initialises process-level C++ state that is not fork-safe. Structured JSON logging via \lstinline|python-json-logger| is configured before any routers are imported.

\subsection{Router Structure}

Nine routers handle all HTTP endpoints. Table~\ref{tab:routers} summarises each router's prefix, purpose, and authentication status.

\begin{table}[H]
  \centering
  \caption{Backend API routers.}
  \label{tab:routers}
  \rowcolors{2}{gray!8}{white}
  \begin{tabularx}{\textwidth}{llX}
    \toprule
    \textbf{Prefix} & \textbf{Auth} & \textbf{Purpose} \\
    \midrule
    \lstinline|/mvp| & No & Video analysis: upload, poll, artefact download \\
    \lstinline|/api/analysis| & \textbf{Yes} & Commit analysis result to Supabase \\
    \lstinline|/api/chat| & \textbf{Yes} & Streaming LLM coaching chat (SSE) \\
    \lstinline|/api/history| & No* & Paginated analysis history retrieval \\
    \lstinline|/api/user| & No* & Profile, stats, account deletion \\
    \lstinline|/api/feedback| & No & Per-analysis feedback retrieval \\
    \lstinline|/api/sessions| & No & Session metadata \\
    \lstinline|/api/recommend| & No & Drill recommendations \\
    \lstinline|/health| & No & Liveness check, version, Gemini status \\
    \bottomrule
    \multicolumn{3}{l}{\small * Identified as an architectural gap; remediation documented in Chapter~\ref{chap:conclusion}.}
  \end{tabularx}
\end{table}

\subsection{Concurrency Model}

Video analysis is CPU-intensive and runs inside a \lstinline|ProcessPoolExecutor| with \lstinline|max_workers = cpu_count - 1|. Jobs are submitted from within an asyncio \lstinline|BackgroundTask|. An asyncio \lstinline|Semaphore(8)| gates submissions, returning HTTP~429 immediately when all slots are occupied rather than queueing, to keep latency predictable.

\subsection{Job Lifecycle}

Figure~\ref{fig:job_lifecycle} illustrates the full job lifecycle from video submission to Supabase persistence.

\begin{figure}[H]
  \centering
  \figurebox[0.83\textwidth]{images/job_lifecycle.png}
  \caption{Analysis job lifecycle: from upload to Supabase persistence.}
  \label{fig:job_lifecycle}
\end{figure}

The numbered steps are:
\begin{enumerate}
  \item \lstinline|POST /mvp/analyze| validates the video, generates a \lstinline|job_id|, inserts \lstinline|status=queued| into SQLite, and returns immediately.
  \item The background task calls \lstinline|MVPPipeline.process_video()| in the process pool (blocking, 15--40\,s typical).
  \item On success, the job service calls Gemini~2.5~Flash with a 10-second timeout to enrich the result, then updates SQLite to \lstinline|status=completed|.
  \item The mobile client polls \lstinline|GET /mvp/result/{job_id}| every two seconds until completion.
  \item After receiving results, the authenticated client calls \lstinline|POST /api/analysis/complete|, which commits the full result to Supabase's \lstinline|analysis_summaries| table and updates the user's streak.
\end{enumerate}

\section{Database and Storage Design}
\label{sec:database}

\subsection{Schema Overview}

The primary database is PostgreSQL managed by Supabase, with Row-Level Security (RLS) enabled on all user-scoped tables. Table~\ref{tab:db_schema} lists the core tables and their purposes.

\begin{longtable}{lp{0.65\linewidth}}
  \caption{Core database tables.}
  \label{tab:db_schema} \\
  \toprule
  \textbf{Table} & \textbf{Purpose} \\
  \midrule
  \endfirsthead
  \toprule
  \textbf{Table} & \textbf{Purpose} \\
  \midrule
  \endhead
  \bottomrule
  \endfoot
  \lstinline|public.users| & SHOOTRZ user records, linked to \lstinline|auth.users| via trigger \\
  \lstinline|user_profiles| & Extended profile (height, weight, experience, coaching style) \\
  \lstinline|user_preferences| & Display and notification settings \\
  \lstinline|user_streaks| & Daily/weekly analysis counters and streak tracking \\
  \lstinline|sessions| & Per-analysis session metadata \\
  \lstinline|videos| & Uploaded video references and Supabase Storage URLs \\
  \lstinline|analysis_summaries| & Full pipeline output: scores, metrics JSON, feedback JSON \\
  \lstinline|chat_history| & Conversational coaching messages (user and coach sides) \\
  \lstinline|drill_library| & Static drill catalogue (read-only for all authenticated users) \\
  \lstinline|drill_completions| & User drill completion history \\
  \lstinline|workout_sessions| & User-defined workout records \\
\end{longtable}

\subsection{Row-Level Security}

Every user-scoped table enforces \lstinline|auth.uid() = user_id| for \lstinline|SELECT|, \lstinline|INSERT|, \lstinline|UPDATE|, and \lstinline|DELETE|. The backend writes through the service-role client (which bypasses RLS), while the frontend anon client is used only for JWT verification.

\subsection{RPC Functions}

Three PostgreSQL functions are defined to reduce round trips:

\begin{itemize}
  \item \lstinline|get_coach_context(user_id, max_analyses)|: Returns recent analyses with metrics for chat context injection, replacing four sequential API calls with a single round trip.
  \item \lstinline|get_user_stats(user_id)|: Returns aggregated streak and analysis statistics.
  \item \lstinline|update_user_streak(user_id, session_id, score)|: Atomically updates streak counters.
\end{itemize}

\subsection{Ephemeral Job Store}

SQLite (via \lstinline|backend/services/job_store.py|) provides a thread-safe, durable in-process store for the lifecycle of analysis jobs with a 72-hour TTL. This avoids the overhead of a full database for ephemeral state while surviving server restarts within the retention window.

\section{AI and ML System Design}
\label{sec:ai_design}

\subsection{MediaPipe Pose Estimation}

MediaPipe Pose (version 0.10.14) provides the pose backbone of the pipeline. It detects 33 full-body landmarks per frame, each with 2D image coordinates, inferred depth, and a visibility confidence score in $[0, 1]$. The system runs in static image mode with model complexity set to 1 (balanced) in production.

Frames with mean landmark visibility below 0.40 are discarded before downstream processing. The shooting side (left or right) is automatically detected from shoulder-to-wrist asymmetry and confirmed against the user's profile. All subsequent angle computations and shot detection prioritise the \lstinline|{shooting_side}_{joint}| landmark naming convention.

\subsection{Shot Event Detection}

Rather than relying on a fixed frame offset, SHOOTRZ implements a multi-signal temporal fusion state machine with four states:

\[
  \text{STANCE} \;\rightarrow\; \text{CROUCH} \;\rightarrow\; \text{RELEASE} \;\rightarrow\; \text{LANDING}
\]

The transition signals are:
\begin{itemize}
  \item \textbf{Knee flexion} (primary crouch signal): knee angle drops below an adaptive threshold derived from the video's joint angle distribution.
  \item \textbf{Hip vertical descent}: hip landmark drops during the loading phase.
  \item \textbf{Wrist height trajectory}: peak wrist height is detected using \lstinline|scipy.signal.find_peaks| over the wrist $y$-coordinate series.
  \item \textbf{Elbow extension}: ascent of elbow angle after the crouch minimum confirms the release phase.
\end{itemize}

Adaptive thresholds are derived per video from signal statistics rather than hardcoded, making the detector resilient to different recording distances and player body proportions.

\subsection{Joint Angle Computation}

Angles are computed from triplets of landmarks using the dot-product formula. For joints A, B, C (with the angle measured at B):

\begin{equation}
  \theta = \arccos\!\left(\frac{(\mathbf{A}-\mathbf{B}) \cdot (\mathbf{C}-\mathbf{B})}{\|\mathbf{A}-\mathbf{B}\|\;\|\mathbf{C}-\mathbf{B}\|}\right)
  \label{eq:angle}
\end{equation}

The argument is clipped to $[-1, 1]$ to prevent numerical errors from \lstinline|arccos| domain violations. Table~\ref{tab:normative} lists the five primary metrics with their optimal ranges and literature sources.

\begin{table}[H]
  \centering
  \caption{Primary biomechanical metrics with optimal ranges.}
  \label{tab:normative}
  \rowcolors{2}{gray!8}{white}
  \begin{tabularx}{\textwidth}{l>{\raggedright\arraybackslash}Xcc}
    \toprule
    \textbf{Metric} & \textbf{Landmarks} & \textbf{Optimal Range} & \textbf{Source} \\
    \midrule
    Elbow extension at release & shoulder $\rightarrow$ elbow $\rightarrow$ wrist & 170--180\textdegree & \cite{Cabarkapa2021} \\
    Knee flexion at crouch & hip $\rightarrow$ knee $\rightarrow$ ankle & 105--115\textdegree & \cite{Cabarkapa2021} \\
    Wrist follow-through & wrist angle change (release $\rightarrow$ landing) & 18--28\textdegree & \cite{Struzik2014} \\
    Forearm verticality & elbow $\rightarrow$ wrist vs.\ vertical & 0--8\textdegree & \cite{Cabarkapa2021} \\
    Release angle (optional) & ball trajectory from YOLOv8n & 63--68\textdegree{} at 4.6\,m & \cite{Okazaki2012} \\
    \bottomrule
  \end{tabularx}
\end{table}

\subsection{Metric Scoring}

Scoring converts raw joint angles into component scores in $[0, 100]$ using a softened distance function:

\begin{enumerate}
  \item If the angle falls within the \lstinline|optimal_range|, score $= 100$.
  \item If within the \lstinline|good_range| but outside optimal, score is linearly interpolated from 100 to 60.
  \item Beyond \lstinline|good_range|, score is penalised proportionally to deviation.
\end{enumerate}

The overall shot score is the confidence-weighted geometric mean of all available component scores. Metric weights are: elbow extension 0.32, knee bend 0.24, release angle 0.20, wrist follow-through 0.14, forearm verticality 0.10. Weights are re-normalised over the subset of metrics with a valid selected frame, so disabling ball tracking does not collapse the score.

Normative ranges are tiered by user skill level. For knee flexion, for example:

\begin{table}[H]
  \centering
  \caption{Skill-tiered normative targets for knee flexion.}
  \label{tab:skill_tiers}
  \rowcolors{2}{gray!8}{white}
  \begin{tabular}{lc}
    \toprule
    \textbf{Skill Level} & \textbf{Target Range} \\
    \midrule
    Beginner & 78--142\textdegree \\
    Intermediate & 89--131\textdegree \\
    Advanced & 100--120\textdegree \\
    \bottomrule
  \end{tabular}
\end{table}

\subsection{LLM Integration}

Gemini 2.5 Flash~\cite{Gemini2024} serves two purposes: generating shot feedback and providing a holistic score that corrects the conservatism of the geometric mean. The client (\lstinline|backend/services/llm/gemini_client.py|) uses the official \lstinline|google-genai| SDK with structured output constrained to two Pydantic schemas:

\begin{itemize}
  \item \lstinline|ShotFeedbackOutput|: overall explanation, metric explanations, strengths list, improvements list, feedback bullets, score tier, AI overall score (0--100), AI score rationale.
  \item \lstinline|DrillRecommendationOutput|: drill name, rationale, step-by-step instructions.
\end{itemize}

All LLM requests are subject to a 10-second timeout with up to three exponential-backoff retries. On failure, the API response sets \lstinline|ai_scored=false| and the rule-based score is used.

\subsection{Drill Recommendation System}

The recommendation engine combines FAISS similarity search~\cite{Johnson2019} with LinUCB contextual bandit reinforcement learning~\cite{Li2010} via the MABWiser library.

A flat inner-product FAISS index (\lstinline|IndexFlatIP|) is built over L2-normalised 5-dimensional user embeddings: \lstinline|[elbow_strength, knee_power, wrist_control, accuracy, consistency]|. The index is queried for the 10 most similar drills from a curated pool of approximately 50 drills.

LinUCB arms correspond to \lstinline|C{cluster}_T{tier}| combinations (drill cluster $\times$ difficulty tier), with exploration parameter $\alpha = 1.4$. The recommendation flow is:

\begin{enumerate}
  \item Normalise the user's metric vector.
  \item FAISS search returns the 10 most similar drills.
  \item LinUCB selects the best-expected arm for the current context.
  \item Candidates are filtered to the chosen arm's cluster and tier.
  \item The top-ranked drill is returned with a natural-language rationale.
\end{enumerate}

\section{Design Decisions and Tradeoffs}
\label{sec:design_decisions}

\subsection{Asynchronous Polling vs.\ WebSocket}

The current design uses HTTP polling (every 2 seconds) rather than WebSocket or SSE for job completion notification. Polling was chosen for simplicity: it requires no persistent connection management, works through all proxies, and is robust to client network interruptions. The trade-off is slightly higher server load from repeated GET requests and a variable 0--2 second notification latency.

\subsection{SQLite as the Ephemeral Job Store}

SQLite provides a durable, thread-safe, zero-dependency storage layer for in-flight jobs that survives backend restarts within the 72-hour retention window and requires no additional infrastructure. The trade-off is that it cannot be shared across horizontally-scaled backend instances, limiting the current architecture to single-instance deployment.

\subsection{Geometric Mean with AI Override}

The geometric mean was chosen over a simple weighted sum because it naturally penalises any single very weak metric---a desirable property for mechanical assessment where a gross form error in one joint genuinely degrades shot quality. However, the geometric mean can be overly conservative when pose estimation confidence is low on one axis but the player's actual form is good. Gemini's holistic score provides a correction for these edge cases, while the rule-based score is always preserved as a deterministic fallback.

\subsection{MediaPipe over Other Pose Estimators}

MediaPipe Pose was selected over alternatives (HRNet~\cite{Sun2019}, OpenPose~\cite{Cao2017}) for its combination of real-time performance on CPU-only hardware, cross-platform support, and well-documented 33-landmark model. The trade-off is that its depth estimation is inferred rather than stereo-captured, limiting 3D measurement accuracy. HRNet is present in the repository as a placeholder for a potential future upgrade.

\subsection{FAISS + LinUCB over Rule-Based Recommendation}

A simple rule-based drill assignment (e.g., ``if elbow $< $ threshold, recommend elbow drill'') cannot adapt to the joint distribution of weaknesses or learn from user engagement over time. The FAISS + LinUCB combination allows recommendations to improve as usage data accumulates, at the cost of additional infrastructure complexity.
